\chapter*{Introduzione}

Il presente lavoro affronta il problema del dimensionamento preliminare del sistema propulsivo di un velivolo civile supersonico. La configurazione considerata prevede un motore ibrido turbogetto--ramjet sviluppato per soddisfare i requisiti di una missione caratterizzata da due differenti condizioni operative. In modo particolare il sistema deve garantire una spinta pari a $25\ \mathrm{kN}$ per motore durante la fase di crociera subsonica a quota $12\ \mathrm{km}$ e numero di Mach $M_0=0.85$. In regime supersonico, si deve garantire una spinta pari a $65\ \mathrm{kN}$ a quota $25\ \mathrm{km}$ e numero di Mach $M_0=3.5$. La missione prevede inoltre una tratta in supercruise di $5000\ \mathrm{km}$ rispettando il limite massimo di combustibile assegnato pari a $24000\ \mathrm{kg}$ per motore.

Per valutare la fattibilità della configurazione proposta sono stati sviluppati modelli dei principali componenti del motore, attraverso i quali sono state determinate le grandezze del ciclo termodinamico e le principali caratteristiche geometriche del motore. 

Il dimensionamento del sistema propulsivo è stato sviluppato mediante una procedura di ottimizzazione vincolata implementata in  MATLAB. A differenza di un approccio  nel quale il ciclo termodinamico viene definito preliminarmente e il dimensionamento dei componenti viene effettuato in una fase successiva, nel presente lavoro le prestazioni del ciclo e la configurazione geometrica del motore sono state determinate simultaneamente all'interno dello stesso algoritmo. Ogni configurazione generata  viene valutata non solo in termini di spinta e consumo di combustibile, ma anche rispetto ai principali vincoli geometrici e aerodinamici dei componenti.


Nei capitoli successivi vengono descritti la metodologia di ottimizzazione adottata, i modelli utilizzati per il dimensionamento dei singoli componenti e i risultati ottenuti per il turbogetto e per il ramjet.